% ===========================================================================
\section{Methodological context}
\label{sec:related}
% ===========================================================================

\subsection{From broad ranges to segment and practice factors}

The 2006 method was defensible when production was the only widely available activity, but its
developing-country ranges combine differences in source coverage, technology and practice. A
geometric midpoint can represent the centre of a bounded log interval; it cannot identify which
physical source is responsible for the result.

The 2019 Refinement changes that. For onshore gas production, for example, the methane factor is
\num{4.09} tonnes CH$_4$ per million m$^3$ for higher-emitting practices and
\SI{2.54}{} for lower-emitting practices, both with reported uncertainty of
$\pm\SI{20}{\percent}$. Distribution can be calculated from gas consumption or pipeline length
and is divided by plastic-pipeline and LDAR conditions. Transmission and processing similarly
distinguish LDAR and compressor-seal practice. Where these characteristics are unknown, the
guidance supports the higher-emitting/default selection and requires compilers to document it
\citep{ipcc2019}.

\subsection{Observation and inventory answer different questions}

Bottom-up and atmospheric estimates often disagree because average factors omit abnormal or
intermittent sources and because measurements sample different source populations. Aircraft and
facility synthesis in the United States found supply-chain methane materially above the national
inventory \citep{alvarez2018methane}; other work showed that a small, heavy-tailed source
population can reconcile apparently divergent estimates \citep{zavalaaraiza2015}. Large sources
also vary in time, so persistence cannot be inferred from a single overpass
\citep{cusworth2021intermittency}.

Satellite observations have made these releases independently visible. TROPOMI provides regional
methane sensitivity \citep{veefkind2012}; global analyses identify ultra-emitters
\citep{lauvaux2022} and national fossil-fuel corrections \citep{chen2023national}. Such
detections do not automatically produce annual national emissions: coverage, detection
thresholds, intermittency and attribution still matter. They can, however, test whether an
inventory's facility geography and source ranking are plausible. Global gridded fuel inventories
face the complementary problem of spatially allocating national reports
\citep{scarpelli2020,scarpelli2022}.

The same distinction applies to flare observations. The World Bank's Global Gas Flaring Tracker
derives annual flare volumes and individual flare locations from satellite radiance
\citep{worldbankflaring2026}. Because the 2019 production factors are inclusive of leakage,
venting and flaring, adding observation-derived flare emissions to them would double count.
Observed flaring is therefore retained as a memo and evaluation layer unless the inclusive
factor is first disaggregated.

\subsection{Official reporting as an audited comparator}

Pakistan's Biennial Transparency Report and CRT provide the authoritative national account
\citep{pakbtr2024,pakcrt2026}. Authority does not make a comparison scope-matched automatically.
The activity units, factors, notation keys, inventory year and GWP basis must be extracted and
compared explicitly. We treat the official submission as a comparator whose derivation is to be
reconstructed, not as a target to which the bottom-up model is forced.
